% !PRINTED 278
\weberpage{278}
Unter dieser Voraussetzung führen die Formeln §. 57 zur
Aufstellung einer Octaëdergruppe.

Die erzeugenden Substitutionen sind:
\[
\Theta = \begin{pmatrix} \gamma^{\frac{p^2-1}{8}}, & 0 \\ 0, & \gamma^{-\frac{p^2-1}{8}} \end{pmatrix}, \quad
\psi = \begin{pmatrix} 0, & \gamma^{\frac{p^2-1}{4}} \\ \gamma^{\frac{p^2-1}{4}}, & 0 \end{pmatrix}
\]
\weq{8}{
\chi = \begin{pmatrix} \dfrac{1}{\sqrt{2}}\,\gamma^{-\frac{p^2-1}{8}}, & \dfrac{1}{\sqrt{2}}\,\gamma^{-\frac{p^2-1}{8}} \\[2mm] \dfrac{-1}{\sqrt{2}}\,\gamma^{\frac{p^2-1}{8}}, & \dfrac{1}{\sqrt{2}}\,\gamma^{\frac{p^2-1}{8}} \end{pmatrix}.
}
Diese Gruppe ist reell, wenn $p\equiv 1\ (\mathrm{mod}\ 8)$ ist, und imaginär, aber in $\Gamma_p$ enthalten, wenn $p\equiv -1\ (\mathrm{mod}\ 8)$ ist.

Das erste Beispiel ist $p=7$. Hier ist $-1$ quadratischer Nichtrest, und man kann also $\varepsilon=\sqrt{-1}=i$ setzen. Nehmen wir $\gamma=2+i$, $\gamma^7=2-i$, $\gamma^6=2\,(1+i)$, $\gamma^8=5$, so ist $\gamma$ eine primitive Wurzel, da $\gamma,\gamma^2,\dots,\gamma^7$ von einander verschieden sind und $5$ reelle primitive Wurzel von $7$ ist, so dass jede nicht verschwindende Zahl des Körpers $\mathfrak{E}_{2,7}$ in der Form
\[
5^{\mu}\gamma^{\nu} = \gamma^{8\mu+\nu}
\]
dargestellt werden kann, wenn
\[
\begin{gathered}
\mu = 0,\ 1,\ 2,\ 3,\ 4,\ 5,\\
\nu = 0,\ 1,\ 2,\ 3,\ 4,\ 5,\ 6,\ 7
\end{gathered}
\]
gesetzt wird.

Es ist dann ferner $\sqrt{2}=3$, und folglich sind die Substitutionen $\Theta,\psi,\chi$:
\weq{9}{
\begin{pmatrix} 2\,(1+i), & 0 \\ 0, & 2\,(1-i) \end{pmatrix}, \
\begin{pmatrix} 0, & i \\ i, & 0 \end{pmatrix}, \
\begin{pmatrix} 4\,(1-i), & 4\,(1-i) \\ -4\,(1+i), & 4\,(1+i) \end{pmatrix}.
}
Will man zur reellen Form übergehen, so muss man nach §. 69 verfahren. Man geht von den Substitutionen $B$ zu den $A$ über, indem man $\nu=\gamma^3=2-3i$ setzt, und erhält aus (9) für $\Theta,\psi,\chi$ in der Gruppe $\Lambda_p$:
\weq{10}{
\begin{pmatrix} 2+2i, & 0 \\ 0, & 2-2i \end{pmatrix}, \
\begin{pmatrix} 0, & 3-2i \\ 3+2i, & 0 \end{pmatrix}, \
\begin{pmatrix} 4-4i, & 1-4i \\ 1+4i, & 4+4i \end{pmatrix}.
}
Hier ist ferner
\[
R = \begin{pmatrix} 4, & 4i \\ i, & 1 \end{pmatrix}
\]
%CONT>>
